\section{Proposed Framework} \label{sec:03-framework}

\begin{figure}[ht] % fig:03-framework-teaser
  \centering
  \includegraphics[width=0.99\textwidth]{./figures/03-framework-teaser.pdf}
  \caption{
    Illustration of \textbf{DIP-KD} of three phases.
    (\textbf{1})~\textbf{Synthesis}:{
      We craft image priors $\mathcal{X}_\mathcal{P}$ with three sub-components: hierarchical noise, nonlinear transformation, and semantic cutmixing.
      Synthetic images $x \sim \mathcal{X}_\mathcal{P}$ have diverse and distinct patterns.
    }
    (\textbf{2})~\textbf{Contrast}:{
      We construct the dataset $\mathcal{D} = \{x | x \sim \mathcal{X}_\mathcal{P}\}$ and train a primer student $S_0$ as a feature extractor.
      Then, we enforce an instance-discriminator $R$ to distinguish embeddings of images $x \sim \mathcal{D}$ to be similar to its positive view $x^{+}$, and dissimilar to its negative view $x^{-}\ne x$.
      Thus, $\mathcal{D}$ is optimized to be even more diverse.
    }
    (\textbf{3})~\textbf{Distillation}:{
      We distill the student $S$ by matching with both hard labels from teacher $T$ via $\mathcal{L}_\text{Hard-KD}$ and soft labels from primer student $S_0$ via $\mathcal{L}_\text{Soft-KD}$, providing extra knowledge.
    }
  }
  \label{fig:03-framework-teaser}
\end{figure}

\subsection{Problem setup} \label{sec:03-framework-setup}
\begin{foldable} % Problem statement
  \textbf{Problem statement.}
  Given solely a \textit{black-box} pre-trained teacher network $T$ that returns only top-1 hard labels (no internal features, logits, or gradients) and \textit{no real training data}, our goal is to train a student $S$ to approximate $T$.
\end{foldable}

\begin{foldable} % Naive solution
  \textbf{Naive solution.}
  Lacking a direct solution, one might adapt standard KD~\cite{hinton2015distilling} as a naive baseline.
  Without real data, this necessitates using uniform noise images, \eg, $x \sim \mathcal{X}_\mathcal{U}$, where $\mathcal{X}_\mathcal{U} = \mathcal{U}[0,1]^{C\times H\times W}$.
  In this restricted interface, only the hard label $\hat{y}_{T}^\text{hard} = T(x) \in \{1, \ldots, K\}$ is available to guide the student's probabilistic outputs $\hat{y}_{S}^\text{soft} = S(x)$.
  The objective is defined as:
  \begin{equation} \label{eq:loss-naive-kd}
    \mathcal{L}_\text{NaiveKD} = \mathbb{E}_{x \sim \mathcal{X}_\mathcal{U}} \left[\mathcal{L}_\text{CE} \left(\hat{y}_{S}^\text{soft}, \hat{y}_{T}^\text{hard} \right)\right],
  \end{equation}
  where $\mathcal{L}_\text{CE}$ is the cross-entropy loss.
  While this can handle trivial cases, it apparently utilizes semantically non-diverse noise, and fails to leverage inter-class knowledge---the primary advantage of KD.
\end{foldable}

\begin{foldable} % Proposal
  \textbf{Proposal.}
  We propose \textbf{Diverse Image Priors Knowledge Distillation} (\textbf{DIP-KD}) to tackle the black-box data-free KD challenge, structured into three collaborative phases:
  \begin{enumerate}
    \item {
        In the \textbf{Synthesis} phase, we design a pipeline to create \textit{image priors} $\mathcal{X}_\mathcal{P}$ that simulate the hierarchical structures, nonlinearity, and meaningful semantics in natural scenes and objects.
        This is achieved via three sub-components: hierarchical noise, nonlinear transformation, and semantic cutmixing.
      }
    \item {
        Then, in the \textbf{Contrast} phase, we sample a dataset $\mathcal{D} = \{x | x \sim \mathcal{X}_\mathcal{P}\}$.
        We further train a \textit{primer student} $S_0$ to optimize $\mathcal{D}$ via a contrastive learning objective to make every sample informationally \textit{distinguishable}.
      }
    \item {
        Finally, in the \textbf{Distillation} phase, we utilize the optimized dataset $\mathcal{D}$ to distill the student $S$.
        Crucially, we extract \textit{soft} knowledge from the primer student $S_0$ to restore the inter-class signals absent in previous BBDFKD methodologies.
      }
  \end{enumerate}
\end{foldable}

\subsection{Synthesis: Reconstructing the knowledge domain} \label{sec:03-framework-synthesis}
\begin{figure}[ht] % fig:03-framework-viz-synthesis
  \centering
  \includegraphics[width=0.99\textwidth]{./figures/03-framework-viz-synthesis.pdf}
  \caption{
    In the \textbf{Synthesis} phase, we generate image priors $\mathcal{X}_\mathcal{P}$ with 3 sub-components.
    (\textbf{1})~\textbf{Hierarchical noise}: We sample $x_\text{hn}$ from multi-scale noise images to capture hierarchical structures.
    (\textbf{2})~\textbf{Nonlinear transform}: We transform $x_\text{hn}$ to $x_\text{nt}$ with nonlinear filters to capture nonlinearity.
    and (\textbf{3})~\textbf{Semantic cutmixing}: We apply cutmixing on $x_\text{nt}$ with a semantic mask $\hat{m}$ to construct the final images $x$, simulating diverse semantics.
    Our image priors $\mathcal{X}_\mathcal{P}$ are visually much more diverse than noise images sampled from random Gaussian or random uniform distributions $\mathcal{X}_{\mathcal{N}}, \mathcal{X}_\mathcal{U}$.
  }
  \label{fig:03-framework-viz-synthesis}
\end{figure}

\begin{foldable} % Motivation --> overview of 3 components
  Natural images are characterized by strong local pixel structures and complex, nonlinear semantics~\cite{lecun2002gradient}.
  To capture these properties, we generate image priors $\mathcal{X}_\mathcal{P}$ through a pipeline of three stages: \textit{hierarchical noise}, \textit{nonlinear transformation}, and \textit{semantic cutmixing}, visualized in Fig.~\ref{fig:03-framework-viz-synthesis}.
  This pipeline serves to capture the structural and semantic breadth of the expert's original knowledge domain through enhanced diversity.
\end{foldable}

\subsubsection{Hierarchical noise} \label{sec:03-framework-synthesis-hiernoise}
\begin{foldable} % Hierarchical structure in nature
  Hierarchy is a universal trait in natural scenes and objects.
  Based on this observation, we design a sampling technique that captures image patterns at both local and global scales to mimic the \textit{hierarchical structures} of visual signals.
\end{foldable}

\begin{foldable} % Sampling multi-scale noise images
  In a standard pixel-wise noise image $x \sim X_\mathcal{U}$, patterns are primarily local and independent.
  Conversely, a monochromatic (single-color) image has dominant global correlation.
  To balance these patterns, we design a sampler that combines noise at multiple scales.
  For an image of dimension $C \times H \times W$, we sample noise images with increasing scales:
  \begin{equation} \label{eq:synthesis-hiernoise-multiscale}
    \left\{ x_{d} |x_{d} \sim \mathcal{U}[0,1]^{C\times2^{d}\times2^{d}} \right\}_{d=0}^{d_\text{max}},
  \end{equation}
  where $d_\text{max} = \text{ceil}(\log_{2} \max(H, W))$ is the first scale to sufficiently cover $H\times W$.
  The dimensions are derived deterministically, \eg, a $32 \times 32$ image is generated from scales $\{2^0, 2^1, \ldots, 2^5\}$.
  To combine local-to-global noise into a single image, we sample coefficients $\alpha_d$ from a normal distribution and apply softmax to get weights $\bar{\alpha}_{d}$.
  We then mix the noise images $\{x_d\}$ to obtain the \textit{hierarchical noise} image:
  \begin{equation} \label{eq:synthesis-hiernoise-combine}
    x_\text{hn} = \sum_{d=0}^{d_\text{max}} \bar{\alpha}_{d} \cdot f_\text{upscale}\left(x_{d};2^{d_\text{max}}\right),
  \end{equation}
  where $f_\text{upscale}$ upscales all noise images to the same dimension $2^{d_\text{max}}$ for proper mixing.
  This approach provides a multi-scale structural foundation that allows the student to learn complex visual hierarchies in the absence of real-world data.
\end{foldable}

\subsubsection{Nonlinear transformation} \label{sec:03-framework-synthesis-nonlintrans}
\begin{foldable} % Nonlinear transformation
  To enrich nonlinear patterns of synthetic images, we apply random rotation $f_\text{rot}$ of $[-45^{\circ}, 45^{\circ}]$, random elastic transform $f_\text{elas}$, and random cropping $f_\text{crop}$ to the specified size $H \times W$.
  These operations are known to be nonlinear and effective in data augmentation~\cite{simard2003best}, and in our case, boosting data diversity.
  Accordingly, we create a \textit{nonlinear-transformed} image as:
  \begin{equation} \label{eq:synthesis-nonlintrans}
    x_\text{nt} = f_\text{crop}^{H\times W} \circ f_\text{elas} \circ f_\text{rot} (x_\text{hn}).
  \end{equation}
\end{foldable}

\subsubsection{Semantic cutmixing} \label{sec:03-framework-synthesis-semcutmix}
\begin{foldable} % Motivation
  We aim to augment synthetic images with the cohesive structures often found in natural objects.
  However, previous techniques like CutMix~\cite{yun2019cutmix} rely on simple rectangular occlusions that lack semantic meaning.
  To address this, we extend CutMix by improving the masking process to nonlinearly blend image pairs into \textit{semantically coherent} structures.
\end{foldable}

\begin{foldable} % Design: diverging filter, mask, and cutmix
  We achieve these patterns through refinement of a semantic mask $\hat{m}$, initialized as $m_{0} \sim \mathcal{U}[0,1]^{1 \times H \times W}$.
  To balance spatial continuity with regional partitioning, we respectively use a blurring filter $f_\text{blur}$ and our \textit{diverging filter} $f_\text{divg}$.
  Specifically, $f_\text{divg}$ is a piecewise quadratic operator that is continuously differentiable at its inflection point (0.5), driving pixel values toward binary states (0 or 1), producing sharp, well-defined region boundaries.
  \begin{equation} \label{eq:synthesis-diverging-filter}
    f_\text{divg}(m) =
    \begin{cases}
      2 m^2           & , m \ge 0.5 \\
      1 - 2 (m - 1)^2 & , m < 0.5
    \end{cases}.
  \end{equation}
\end{foldable}

\begin{foldable} % Steady-state convergence and Diversity
  We obtain the final $\hat{m}$ by iteratively applying $m_{i+1} = f_\text{divg} \circ f_\text{blur}(m_i)$ until the mean absolute pixel-wise update reaches a steady-state threshold $\epsilon < 10^{-2}$, empirically around 10 iterations.
  This mask is used for pairwise cutmixing $\{x_\text{nt}\}$ against themselves or random monochromatic images $\mathcal{X}_\text{mc}$ to form our final \textit{image priors} $\mathcal{X}_\mathcal{P} = \{x\}$, where:
  \begin{equation} \label{eq:synthesis-semcutmix}
    x= \text{CutMix}(x_i, x_j; \hat{m}), \text{\quad s.t. \quad} x_i,x_j \sim \{x_\text{nt}\} \cup \mathcal{X}_\text{mc} \text{ and } x_i \ne x_j.
  \end{equation}
  This approach ensures the synthetic domain of $\mathcal{X}_\mathcal{P}$ captures diverse semantics required for effective black-box knowledge acquisition.
\end{foldable}

\subsection{Contrast: Collaborative diversity optimization} \label{sec:03-framework-contrast}
\begin{foldable} % Primer student S_0
  Standard contrastive learning methods require internal feature extraction, which is prohibited by the black-box nature of our teacher.
  We propose the novel use of a \textit{primer student} $S_{0}$ as a white-box mediator to bridge this informational gap.
  We first train $S_{0}$ on an image priors dataset $\mathcal{D} = \{x | x \sim \mathcal{X}_\mathcal{P}\}$ following the Hard-KD objective:
  \begin{equation} \label{eq:loss-hard-kd}
    \mathcal{L}_\text{Hard-KD} = \mathbb{E}_{x \sim \mathcal{D}} \left[\mathcal{L}_\text{CE} \left(\hat{y}_{S_0}^\text{soft}, \hat{y}_{T}^\text{hard} \right)\right],
  \end{equation}
  where $\hat{y}_{S_0}^\text{soft} = S_0(x)$ is its probability prediction, and $\hat{y}_{T}^\text{hard} = T(x)$ is the teacher's hard label.
  Equipped with this white-box proxy, we can access its internal representations for explicit diversity optimization.
\end{foldable}

\begin{foldable} % Embeddings & cosine similarity
  Following~\cite{fang2021contrastive}, diversity correlates with how \textit{``distinguishable''} samples are, under an \textit{explicitly optimizable} objective.
  To maximize diversity, we use the primer student's backbone $S_{0}^\mathcal{H}$ to extract image embeddings and append an instance-discriminator $R$ (a shallow MLP) to recognize each embedding as a unique class.
  We argue that even if $S_0$ is yet to match the teacher's expertise, its backbone $S_{0}^\mathcal{H}$ is capable of capturing the underlying semantics required for contrastive learning.
  This is supported by prior works in contrastive learning~\cite{chen2020simple,he2020momentum,zbontar2021barlow}, which do not depend on a high-accuracy classifier, but often train from scratch.
  For any two images $(x_i, x_j)$, we compute their cosine similarity as:
  \begin{equation}
    \text{Sim}\left(x_i, x_j\right)=
    \frac{ \left\langle R \circ S_{0}^\mathcal{H}(x_i), R \circ S_{0}^\mathcal{H}(x_j)\right\rangle }
    { \left\Vert R \circ S_{0}^\mathcal{H}(x_i) \right\Vert \cdot \left\Vert R \circ S_{0}^\mathcal{H}(x_j) \right\Vert }.
  \end{equation}
\end{foldable}

\begin{foldable} % Contrastive loss
  For diversity optimization, from an image $x \sim \mathcal{D}$, we generate positive views $x^{+}$ via augmentation and sample other images $x^{-} \sim \mathcal{D}, x^{-} \neq x$ as negative views.
  The contrastive objective maximizes the similarity between $(x, x^{+})$ and minimizes that between $(x, x^{-})$:
  \begin{equation} \label{eq:loss-contrastive}
    \mathcal{L}_\text{CT} = \mathbb{E}_{x \sim \mathcal{D}} \left[
      -\log \frac{ \exp\left( \text{Sim}(x, x^{+}) \right) }{ \sum_{x^{-} \ne x} \exp\left( \text{Sim}(x, x^{-}) \right) }
    \right].
  \end{equation}
  By backpropagating $\nabla\mathcal{L}_\text{CT}$ to update both the instance-discriminator $R$ and the samples $x \in \mathcal{D}$, we explicitly reinforce the semantic diversity of $\mathcal{D}$.
\end{foldable}

\subsection{Distillation: Restoring informative signals} \label{sec:03-framework-distillation}
\begin{foldable}
  Once the synthetic dataset $\mathcal{D}$ has been optimized for diversity, we use it for the final training of the student $S$.
  Our objective is to restore the high-bandwidth signals typical of standard KD~\cite{hinton2015distilling}, \ie, incorporating both \textit{hard} and \textit{soft} knowledge, which were \textit{limited or entirely absent} in recent works~\cite{wang2021zero, zhang2022ideal, yuan2024data}.
  Having privileged access to the primer network $S_0$, we form an optimal distillation objective by jointly matching the student's logits $S^\mathcal{Z}(x)$ with soft labels $S_{0}^\mathcal{Z}(x)$ and hard labels $\hat{y}_{T}^\text{hard}$:
  \begin{align} \label{eq:loss-kd-total}
    \mathcal{L}_{S}
    &= \mathcal{L}_{\text{Hard-KD}}+\mathcal{L}_{\text{Soft-KD}} \\
    &= \mathbb{E}_{x\sim\mathcal{D}}\Big[\mathcal{L}_{\text{CE}}\left(\hat{y}_{S}^\text{soft},\hat{y}_{T}^\text{hard}\right) + \mathcal{L}_{\text{L1}}\left(S^{\mathcal{Z}}\left(x\right),S_{0}^{\mathcal{Z}}\left(x\right)\right)\Big],
  \end{align}
  where $\mathcal{L}_{\text{L1}}$ is the L1 divergence loss, applied directly to logits to preserve raw inter-class relationships.
  By reconciling the teacher's authoritative expertise with the primer student's relational hints, the student $S$ overcomes the signal scarcity of the black-box interface.
  This collaborative training allows for a deeper understanding of the teacher's decision-making logic, resulting in superior performance across restricted AI ecosystems.
\end{foldable}
